%-------------------------
% CV in Latex (Turkish)
% Author : Bedirhan Yenilmez
% Based off of: https://www.overleaf.com/latex/templates/jakes-resume/syzfjbzwjncs
% License : MIT
%------------------------

\documentclass[a4paper,11pt]{article}

\usepackage[utf8]{inputenc}
\usepackage[T1]{fontenc}
\usepackage{lmodern}

\usepackage{latexsym}
\usepackage[empty]{fullpage}
\usepackage{titlesec}
\usepackage{marvosym}
\usepackage[usenames,dvipsnames]{color}
\usepackage{verbatim}
\usepackage{enumitem}
\usepackage[hidelinks]{hyperref}
\usepackage{fancyhdr}
\usepackage[turkish,shorthands=off]{babel}
\usepackage{tabularx}
\input{glyphtounicode}

\pagestyle{fancy}
\fancyhf{} % clear all header and footer fields
\fancyfoot{}
\renewcommand{\headrulewidth}{0pt}
\renewcommand{\footrulewidth}{0pt}

% Adjust margins
\addtolength{\oddsidemargin}{-0.5in}
\addtolength{\evensidemargin}{-0.5in}
\addtolength{\textwidth}{1in}
\addtolength{\topmargin}{-.65in}
\addtolength{\textheight}{1.3in}

\urlstyle{same}

\raggedbottom
\raggedright
\setlength{\tabcolsep}{0in}
\setlength{\footskip}{5pt}

% Sections formatting
\titleformat{\section}{
  \bfseries\raggedright\large
}{}{0em}{}[\color{black}\titlerule]
% Format: \titlespacing*{<command>}{<left>}{<before-sep>}{<after-sep>}
\titlespacing*{\section}{0pt}{3pt}{4.25pt}

% Ensure that generate pdf is machine readable/ATS parsable
\pdfgentounicode=1

%-------------------------
% Custom commands
\newcommand{\resumeItem}[1]{
  \item\small{#1}
}

\newcommand{\resumeSubheading}[4]{
  \item
    \begin{tabular*}{0.97\textwidth}[t]{l@{\extracolsep{\fill}}r}
      \textbf{#1} & #2 \\
      \textit{\small#3} & \textit{\small #4} \\
    \end{tabular*}
}

\newcommand{\resumeSubSubheading}[2]{
    \item
    \begin{tabular*}{0.97\textwidth}{l@{\extracolsep{\fill}}r}
      \textit{\small#1} & \textit{\small #2} \\
    \end{tabular*}
}

\newcommand{\resumeProjectHeading}[2]{
    \item
    \begin{tabular*}{0.97\textwidth}{l@{\extracolsep{\fill}}r}
      \small#1 & #2 \\
    \end{tabular*}
}

% Control spacing between jobs/projects
\newcommand{\resumeSubHeadingListStart}{
    \begin{itemize}[leftmargin=0.15in, label={}, itemsep=1pt, parsep=0pt, topsep=0pt]
}
\newcommand{\resumeSubHeadingListEnd}{\end{itemize}}

% Control spacing between bullet points under a job/project
\newcommand{\resumeItemListStart}{
    \begin{itemize}[label={$\vcenter{\hbox{\tiny$\bullet$}}$}, itemsep=0.5pt, parsep=0pt, topsep=0.0pt, partopsep=0pt]
}
\newcommand{\resumeItemListEnd}{\end{itemize}}

\linespread{1.0} % Adjusts the spacing between all lines of text

%-------------------------------------------
%%%%%%  RESUME STARTS HERE  %%%%%%%%%%%%%%%%%%%%%%%%%%%%

\begin{document}

%----------HEADING----------
\begin{center}
\href{https://bedirhanyenilmez.com/tr}{\textbf{\Huge \scshape Bedirhan Yenilmez}} \\ \vspace{1pt}
\small
%PHONE%
\href{mailto:yenilmezbedirhan@gmail.com}{\underline{yenilmezbedirhan@gmail.com}} $|$
\href{https://linkedin.com/in/bedirhan-yenilmez}{\underline{linkedin.com/in/bedirhan-yenilmez}} $|$
\href{https://github.com/beyenilmez}{\underline{github.com/beyenilmez}}
\end{center}

%-----------SUMMARY-----------
\section{\smash{Özet}}
\begin{itemize}[leftmargin=0.15in, label={}]
  \small{\item{
    Açık kaynak projelere ve bağımsız geliştirmeye ilgi duyan bir yazılım mühendisi. Kendi projelerimi üretmeyi seviyor; kendi kendime öğrenerek kazandığım bu pratik becerileri ve teknik bilgiyi profesyonel iş hayatımda somut sonuçlara dönüştürmeyi hedefliyorum.
  }}
\end{itemize}

%-----------EDUCATION-----------
\section{\smash{Eğitim}}
  \resumeSubHeadingListStart
    \resumeSubheading
      {Dokuz Eylül Üniversitesi}{İzmir, Türkiye}
      {Bilgisayar Mühendisliği Lisans}{Eyl 2022 -- Haz 2026}
       \resumeItemListStart
        \resumeItem{GNO: 3.8/4.0}
      \resumeItemListEnd
  \resumeSubHeadingListEnd


%-----------EXPERIENCE-----------
\section{Deneyim}
  \resumeSubHeadingListStart

    \resumeSubheading
      {Semafor Teknoloji}{İzmir, Türkiye}
      {Yazılım Mühendisliği Stajyeri}{Ara 2025 -- Şub 2026}
      \resumeItemListStart
        \resumeItem{Saha operasyonlarında otomatik KKD (Kişisel Koruyucu Donanım) uygunluk denetimi için \textbf{YOLOv8} ve \textbf{EasyOCR} kullanarak \textbf{FastAPI} ve \textbf{Gemini LLM} tabanlı bir sesli asistan geliştirdim.}
        \resumeItem{Akıllı saat için \textbf{Kotlin} ve \textbf{Firebase} kullanarak bir sağlık takip sistemi ve \textbf{Next.js} ile etkileşimli bir kontrol paneli oluşturdum.}
      \resumeItemListEnd

    \resumeSubheading
      {ASELSAN}{Ankara, Türkiye}
      {Yazılım Mühendisliği Stajyeri}{Ağu 2025 -- Eyl 2025}
      \resumeItemListStart
        \resumeItem{Özelleştirilmiş veri dönüştürme katmanları kurarak, \textbf{C++} ile modüler ve çift yönlü \textbf{donanım--ROS 2 iletişim arayüzleri} geliştirdim.}
        \resumeItem{Test ve hata ayıklama işlemleri için \textbf{ROS 2 araçlarını} (\textbf{ros2cli}, \textbf{Foxglove}, \textbf{Lichtblick}) kullandım.}
      \resumeItemListEnd

    \resumeSubheading
      {HAVELSAN}{Ankara, Türkiye}
      {Yazılım Mühendisliği Stajyeri}{Tem 2025 -- Ağu 2025}
      \resumeItemListStart
        \resumeItem{\textbf{.NET 9}, \textbf{Avalonia UI} ve \textbf{MVVM} mimarisini kullanarak çapraz platform bir masaüstü uygulaması geliştirdim.}
        \resumeItem{\textbf{Bullseye}, \textbf{Docker} ve \textbf{NSIS} kullanarak derleme ve paketleme süreçlerini otomatikleştirdim.}
      \resumeItemListEnd

  \resumeSubHeadingListEnd


%-----------CERTIFICATIONS-----------
\section{\smash{Sertifikalar}}
  \resumeSubHeadingListStart

    \resumeSubheading
      {\href{https://drdogrulama.sanayi.gov.tr/tr/verify/07249146037237}{Yapay Zeka Uzmanlık Eğitim Sertifikası}}
      {Eyl 2025}
      {T.C. Sanayi ve Teknoloji Bakanlığı}{}
      \resumeItemListStart
        \resumeItem{Yapay zeka uygulamaları, veri yönetimi, dağıtık programlama ve algoritmaları kapsayan \textbf{50 oturumluk Yapay Zeka Uzmanlık Eğitim Programı'nı} tamamladım.}
      \resumeItemListEnd

    \resumeSubheading
      {{Savunma Sanayii 401 Sertifikası}}{Şub 2025}
      {Savunma Sanayii Başkanlığı \& Yükseköğretim Kurulu}{}
      \resumeItemListStart
        \resumeItem{\textbf{Savunma Sanayii 401} programı kapsamında; savunma teknolojileri, sistem mühendisliği, proje yönetimi ve kalite güvencesi alanlarında eğitim aldım.}
      \resumeItemListEnd

  \resumeSubHeadingListEnd


%-----------PROJECTS-----------
\section{\href{https://bedirhanyenilmez.com/tr/\#projects}{\smash{Projeler}}}
    \resumeSubHeadingListStart
      \resumeProjectHeading
          {\textbf{Askeri Kara Araçlarında Yapay Zeka ile Sürücü Yorgunluk Tespiti}}{2026}
          \resumeItemListStart
            \resumeItem{SAYZEK-ATP (Savunma Sanayii Yapay Zeka Yetenek Kümelenmesi -- Akademik Tez Programı) ve TÜBİTAK 2209-A destek programlarına kabul edilen bir \textbf{Lisans Bitirme Projesi} olarak geliştirdim.}
            \resumeItem{Yerel olarak çalışan ve \textbf{0.99}'a varan F1 skoruna ulaşan bir yorgunluk tespit sistemi oluşturdum.}
          \resumeItemListEnd

      \resumeProjectHeading
          {\textbf{\href{https://github.com/beyenilmez/autostartstop}{AutoStartStop}}}{2026}
          \resumeItemListStart
            \resumeItem{Minecraft proxy alt sunucularını dinamik olarak yönetmek için kural tabanlı bir \textbf{Java} eklentisi geliştirdim. \textbf{JUnit} ve \textbf{Mockito} ile kapsamlı birim ve entegrasyon testleri yazdım.}
          \resumeItemListEnd

      \resumeProjectHeading
          {\textbf{\href{https://github.com/beyenilmez/pz-admin}{PZ Admin}}}{2024 -- 2025}
          \resumeItemListStart
            \resumeItem{Project Zomboid sunucularının RCON üzerinden yönetilmesi için çapraz platform bir uygulama geliştirdim. \textbf{GitHub Actions} ile derleme süreçlerini otomatikleştirdim.}
          \resumeItemListEnd

    \resumeSubHeadingListEnd


%-----------SKILLS-----------
\section{\smash{Teknik Beceriler}}
 \begin{itemize}[leftmargin=0.15in, label={}]
    \small{\item{
     \textbf{Diller}{: Go, C\#, Java, Python, C/C++, SQL (PostgreSQL), TypeScript/JavaScript} \\
     \textbf{Frameworkler}{: Avalonia, React, Next.js, Wails} \\
     \textbf{Araçlar \& Teknolojiler}{: Linux, Git, Docker, Proxmox}
    }}
 \end{itemize}

%-----------LANGUAGES-----------
\section{\smash{Diller}}
\begin{itemize}[leftmargin=0.15in, label={}]
  \small{\item{
    \textbf{Türkçe}{: Ana dil} \\
    \textbf{English}{: Orta-İleri Düzey (B2)}
  }}
\end{itemize}

%-----------TECHNICAL INTERESTS-----------
\section{\smash{Teknik İlgi Alanları}}
\begin{itemize}[leftmargin=0.15in, label={}]
  \small{\item{
    \vspace{-0.3mm}Ev sunucusu (Homelab) yönetimi ve açık kaynaklı yazılım katkıları
  }}
\end{itemize}

%-------------------------------------------
\end{document}
